% Abhakseite „Das kann ich“ – Zone nur im Gesamt (\mitzone)
%% TODO Ich-kann-Sätze: Platzhalter wie in den Aufgabendateien
\begin{abhakseite}
\ifdefined\mitzone
\abhakgruppe{Kennst du schon}
\abhak{1}{Koordinaten lesen und eintragen, 4 Quadranten, (x|y)-Reihenfolge}
\abhak{2}{Proportionale Zuordnung erkennen und hochrechnen (Dreisatz)}
\abhak{3}{Negative Zahlen multiplizieren und dividieren}
\abhak{4}{Brüche als Steigung (½, −¾), Bruch mal ganze Zahl}
\abhak{5}{Lineare Gleichung zweischrittig lösen (null gleich einem zweischrittigen Term)}
\abhak{6}{Terme zusammenfassen}
\abhak{7}{Negative Zahlen multiplizieren und dividieren – Fehler finden}
\abhak{8}{Negative Zahlen multiplizieren und dividieren – selbst rechnen}
\fi
\abhakgruppe{Einheit 1 · Proportionale Funktion}
\abhak{9}{Proportionale Funktion}
\abhak{10}{Dreisatz als Kontrolle}
\abhak{11}{Proportionale Funktion – Fehler finden, Begründen, Darstellungswechsel, Anwendung}
\abhakgruppe{Einheit 2 · Lineare Funktion f(x) = m·x + n}
\abhak{12}{m und n in der Gleichung markieren}
\abhak{13}{Steigungsrichtung ohne Zeichnen}
\abhak{14}{Steigungsdreieck lesen}
\abhak{15}{Graph zeichnen}
\abhak{16}{Ablesen}
\abhak{17}{Parameter deuten (steigend, fallend, parallel, Sonderfall m = 0)}
\abhak{18}{Graph zeichnen – Fehler finden, Begründen, Darstellungswechsel}
\abhakgruppe{Einheit 3 · Punkte und Werte}
\abhak{19}{Punkt in Gleichung einsetzen}
\abhak{20}{Funktionswert}
\abhak{21}{Wertetabelle}
\abhak{22}{Schnittpunkt mit y-Achse}
\abhak{23}{Funktionswert – Fehler finden, Begründen, Darstellungswechsel, Anwendung}
\abhakgruppe{Einheit 4 · Gleichung bestimmen}
\abhak{24}{Gleichung bestimmen}
\abhak{25}{Gerade durch zwei Punkte zeichnen}
\abhak{26}{Schnittpunkt zweier Geraden rechnerisch}
\abhak{27}{Gleichung bestimmen – Fehler finden, Begründen, Darstellungswechsel, Anwendung}
\abhakgruppe{Einheit 5 · Anwendungen}
\abhak{28}{Anwendung}
\abhak{29}{Graph zu Tarif zuordnen}
\abhak{30}{Situation zu Gleichung beschreiben}
\abhak{31}{Anwendung – Fehler finden, Begründen, Darstellungswechsel, Anwendung}
\end{abhakseite}
